% vim: tw=78 ai spell spelllang=cs

\documentclass{article}

\usepackage{graphicx}
\usepackage{hyperref}
\usepackage[utf8x]{inputenc}
\usepackage[czech]{babel}
\usepackage[T1]{fontenc}
\usepackage{xspace}

\pagestyle{empty}

\newcommand{\topinfo}[4]{
\begin{flushright}
\today
\end{flushright}
\vspace{1em}
\begin{center}
\Large\textbf{Posudek #1 na #2 práci \\
\normalsize%
\textit{#3:} #4}
\end{center}
}
\newcommand{\botinfo}{
\vspace{3em}
\begin{flushright}
Jiří Slabý
\end{flushright}
}


\newcommand{\cdb}{\textsc{ClabureDB}}

\begin{document}
\topinfo{vedoucího}{bakalářskou}{Viktor Hrtánek}{Podpora pro vkládání záznamů
do \cdb}

Dle zadání měla bakalářská práce nastudovat vkládání dat do databází různých
systémů a doimplementovat takovou funkcionalitu do fakultního nástroje \cdb,
a to s podporou dvou vstupních formátů, přičemž měla být korektně ošetřena
přístupová práva a celá implementace zdokumentována. Text práce je na 21
stranách strukturovaný do 4 kapitol.  V úvodu autor krátce popisuje motivaci
pro práci.  Následující kapitola popisuje různé databáze benchmarků.
Následuje kapitola se stručným popisem implementace a závěr. \textbf{Rozsah}
práce se pohybuje na samé spodní hranici požadavků na bakalářské práce.

\textbf{Obsahově} je dle mého názoru dostatečně propracována pouze kapitola
2, tedy přehledová (i zde ale mohl být rozsáhlejší popis \cdb). Ostatní
kapitoly jsou příliš stručné, jmenovitě chybí dobrý úvod do práce a kvalitní
popis implementace, kde je jen vágně uvedeno, jak implementace funguje, jaký
formát dat je vstupem, co přesněji dělají popsané knihovny (např.\ popis
\textsc{Paperclip} toho příliš nevysvětluje) apod. Několik tvrzení a návrhů
nedává valný smysl:
\begin{itemize}
\item V 3.3: proč se data o uživatelích mají ukládat do jedné z poddatabází --
není to porušením 2.\ NF?
\item V 3.3: na co je tam celý ER diagram, když není téměř vysvětleno, co na
něm je a k čemu to je?
\item Str.\ 23, \textsc{Slim-rails}: \texttt{th} není párovaný.
\item V literatuře: ke \textsc{Stanse} bylo lépe citovat existující článek, ne
odkaz na web.
\item Příloha B: špatné příkazy. Všechny by měly být pomocí \texttt{bundle
exec}. Zejména v bodě 4, \texttt{exec db:seed} vyvolá příkaz \texttt{exec}
příkazového řádku, který skončí chybou.
\end{itemize}

Po \textbf{jazykové} stránce je práce na podprůměrné úrovni. Je zde mnoho
hrubých chyb, první již v 5.\ slově práce. Jako příklad uvádím
,,mechanizmi'', ,,analýzi'', ,,programi'', ,,záznami'' podle vzoru
\emph{dub}, ,,kernely'', ,,adresáry'' podle vzoru \emph{stroj} po řadě na
stranách 3, 11, 11, 16, 16, 17 a 22. Je zde také mnoho překlepů (např.\
,,implentáciou'' na str.\ 1, ,,zoznam požitých gemov'' na str.\ 22,
,,Created'' a ,,Approved'' ,,on'', ne ,,at'' v obr.\ 3.4.1), několik
amerikanismů (,,namespacu'' na str.\ 28).  Občas dochází k řádkům končícím na
jednopísmenné předložky (str.\ 17) a sirotkům (str.\ 24).

\textbf{Praktická část} práce je v přiloženém \texttt{rar} souboru. Přestože
se jedná o rozšíření již existujícího nástroje \cdb, není v práci zcela přesně
uvedeno, co přesně bylo změněno a přidáno. \textsc{Git} napovídá 27 souborů a
579 změněných řádků. Tento nově přidaný kód ale není zcela funkční. Zejména
mám tyto připomínky:
\begin{itemize}
\item \texttt{app/models/project.rb}: k čemu je změna z dynamického načítání
poddatabází směrem k pevně danému přímo v kódu?
\item Při stisknutí tlačítka \emph{Sign out} aplikace havaruje s hlášením:
\begin{verbatim}
ActiveRecord::StatementInvalid (SQLite3::SQLException: no
such table: admins: SELECT  "admins".* FROM "admins"
WHERE "admins"."id" = 5 ORDER BY "admins"."id" ASC LIMIT 1)
\end{verbatim}
\item Heslo se při změně nezmění, jak je psáno v příloze B.
\item Mohu importovat jakýkoliv typ souboru. Pak ho lze stáhnout a prohlížet
(např.\ PDF). Nicméně při stisknutí ,,Approve'' pak aplikace samozřejmě
přestane správně pracovat a objeví se šedá obrazovka s výpisy.
\end{itemize}

Na závěr konstatuji, že přes všechny uvedené připomínky práce \textbf{splňuje
zadání a požadavky na bakalářskou práci}. Zejména proto, že se autor
dobrovolně rozhodl pracovat sám, bez mojí podpory. Doporučuji práci
\textbf{uznat} jako bakalářskou, ale vzhledem k výše uvedeným okolnostem
doporučuji hodnocení stupněm E.

\botinfo
\end{document}
